\section{Introduction}
A portfolio manager rarely knows expected returns with equal precision across assets. A common shrinkage parameter simplifies the response to this uncertainty, but it also removes differences that could affect the allocation. This paper asks how much conditional portfolio utility is lost when heterogeneous signal uncertainty is replaced by the best single adjustment. It then asks whether an estimate of those differences improves net portfolio utility when forecasts, risk inputs, parameter selection and trading costs are held comparable. These questions connect the informational content of uncertainty estimates to their economic value in portfolio decisions.

Mean--variance portfolio selection formalizes the trade-off between expected return and return variance \citep{markowitz1952}. Its implementation requires estimated inputs. Covariance regularization addresses the conditioning of the estimated risk matrix, as developed by \citet{ledoit2004}. Enhanced portfolio optimization (EPO), introduced by \citet{pedersen2021}, relates correlation shrinkage to uncertainty in expected returns and to several existing approaches to portfolio construction. Their analysis identifies portfolios associated with small estimated eigenvalues as important sources of poor conventional optimization performance. In their industry application, the reported gross Sharpe ratio increases from 0.19 for unshrunk mean--variance optimization to 0.79 for EPO with a parameter selected from past data.

The original EPO framework already allows a general covariance matrix for signal errors. Our question concerns the consequences of restricting that matrix to a common scalar times asset variances. The distinction matters because asset risk and precision of a conditional-mean forecast measure different quantities. Two industries with similar return variances can have different uncertainty about their expected returns. Even an optimally chosen common penalty need not preserve the allocation implied by these differences.

We implement three connected stages. First, we reproduce the Equity 1 industry benchmark and its principal-component diagnostics. This stage establishes how shrinkage affects the implemented baseline. Second, in a Gaussian model with known covariance and signal uncertainty, we compare the optimal heterogeneous allocation with the best allocation within the scalar family. The scalar comparator uses the same information as the heterogeneous rule. Its remaining utility loss measures the effect of the scalar restriction; a separate comparison measures the cost of choosing a scalar parameter poorly. Third, we construct a feasible uncertainty estimate from repeated fits of a rolling forecasting regression and compare two industry portfolios that share a single historically selected penalty level.

Our methodological contribution is the combination of an information-matched scalar comparison, a calibration experiment using the actual uncertainty estimator, and a matched economic evaluation after costs. The portfolio objective, EPO allocation, linear regression and block bootstrap are established components. The project specifies how these components are combined to measure scalar compression and to evaluate one implementable heterogeneous rule. In particular, the heterogeneous portfolio receives no separately selected tuning parameter.

The benchmark replication produces gross Sharpe ratios of 0.322 for raw optimization, 0.592 after the preliminary covariance correction, and 0.770 for EPO selected from past data. In the simulations, uniform uncertainty gives essentially zero compression loss. The two heterogeneous profiles produce losses equal to 12.4\% and 28.7\% of mean heterogeneous utility, with larger losses in the high exposure-alignment groups. The feasible comparison produces annual net certainty equivalent differences of 27.65 and $-98.52$ basis points in the two evaluation periods. Both 95\% intervals include zero. The results connect the value of knowing uncertainty differences to the outcomes of estimating and using those differences in a portfolio.
